\subsection{Scope and limitations}

The noncausal model experiments are fixed-input inference evaluations rather
than finetuning or pretraining studies.  The causal measurements establish
operator, complete-sublayer, and full-update performance.  The matched
distributed comparison supports training and validation claims only through
the reporting cutoff in Section~\ref{sec:causal-training}; one trajectory per
route does not estimate run-to-run variability.  The longer historical
controls establish the FP8-versus-MXFP4 P/V format choice, but their different
recipe prevents a paired quality comparison with BF16.  Learned projection
precision is also a separate boundary: E4M3 is the numerical control and
NVFP4 is the measured throughput arm.

The hardware conclusions are shape-specific.  Head dimension 64 and other
regimes require their own tile size, cooperative-thread-array ownership, and
tensor-memory overlap strategy; they should not inherit the D128 schedule by
analogy.

\subsection{Conclusion}

Direct-P makes full-FP4 forward attention useful by shortening the serial path
between the score and value matrix products.  It maps log-softmax scores
directly to MXFP4 codes and normalizes with the same represented probabilities
consumed by the value product.  The result exceeds twice the BF16 forward
throughput on favorable GB200 shapes.  The FP8 probability route remains more
accurate, while the fixed-input model evaluations identify cases in which the
additional Direct-P error remains small at the model output.

For causal training, the forward pass supplies backward with its quantized
Q/K payload, scales, and softmax normalizer.  Projection and gradient
epilogues also publish the row- and column-oriented FP8 views required by the
gradient products.  This path accelerates the complete projection-inclusive
attention sublayer by 1.25$\times$ and a single-GPU 8B update by up to
1.14$\times$.  The distributed controls select FP8 P/V: every tested MXFP4
P/V trajectory diverges despite competitive short timing runs.

Faster matrix arithmetic alone is therefore insufficient.  At head dimension
128, tensor-memory ownership and operand readiness limit how much QK, softmax,
and PV can overlap.  Low-precision attention needs a larger allocatable
overlap window as much as it needs faster tensor cores.
